\section{第五大组组织分工与协作边界}\label{sec:division}
本报告是第五大组共同实验报告，统一汇总TonyPi与各类TurtleBot跟随方式的设计、实现和实验。第一小组负责TonyPi抱物路线与方式二；第二小组方式一的成员与分工采用仓库已提交文稿；第三小组成员信息待补充；方式四小组已有四位成员及学号，行政组号尚待确认。方式编号、小组编号与车辆编号分别登记，避免将三车队列中的“一号车”当作“第一小组”。
\begin{table}[H]\centering\small
\begin{tabularx}{\linewidth}{p{27mm}YY}\toprule
功能模块/责任 & 技术范围 & 协作输入与输出\\\midrule
第一小组：TonyPi & 抱物动作、固定标记识别、路线状态机、转向与放置 & A/B/C标记、动作与日志；向后车提供可识别目标与状态适配\\
第一小组：方式二 & 建图、定位、AR目标转换、轨迹跟随、规划与门控 & 地图、定位、扫描和任务许可；导航速度\\
第二小组：方式一 & AR直接跟随、相机/TF、目标ID与门控方案 & 背贴ALVAR量测；跟随控制、标定与试验资料\\
第三小组：方式三（组号待核） & RGB-D候选、身份关联、滚动地图与控制 & RGB-D和本车里程计；身份目标、规划与状态\\
方式四小组（组号待核） & 点云首车与AR后车组成的三车纵列，扫描响应与有界搜索 & 前车ID0/1、扫描；C++控制器及三车联调说明\\
全队共同 & 场地布置、速度源协调、方式认定、证据汇总与联合验收 & 统一接口表、各车完成记录与共同报告\\\bottomrule
\end{tabularx}\caption{系统功能与模块边界。}
\end{table}
第一小组四位成员姓名由本次委托确认，个人职责参考\code{v1.pdf}封面及第2章（PDF第1、6--7页）归纳；不在缺少依据时指定某位成员独自负责全部地图导航。第二、方式四小组名单来自各组提交，交付记录表示对应组的登记内容，不等于全部原始数据已经归档。第一、第二小组学号待补，方式四学号保留原投稿数值。
各组在独立章节和图片子目录中修改技术内容。公共导言、主文件及引用格式统一维护；跨模块参数变化应同步记录接口影响。队列拓扑、跟随目标来源和导航方法分别登记，避免将同一控制器在不同车号上的运行视为不同方法。

\subsection{第一小组组内分工与协作}
第一小组采用共同设计、交叉调试与协同完善的工作方式，围绕动作标定、二维码转弯、状态机实现、实验环境和报告整理形成相互衔接的任务分工。表中归纳主要承担方向；共同参与的内容不重复计为互相独立的成果，也不作定量贡献比例分配。
\begin{longtable}{p{18mm}p{37mm}p{88mm}}
\caption{第一小组成员的主要分工与交叉协作}\\\toprule
成员 & 主要承担方向 & 工作内容与协作关系\\\midrule\endfirsthead
\toprule 成员 & 主要承担方向 & 工作内容与协作关系\\\midrule\endhead
\textbf{杨开元} & 动作标定、状态机与转弯算法初步实现 & 面向抱物行走与左右转弯开展动作标定，编写路线状态机并初步实现二维码转弯算法；衔接标记观测、阶段判断与动作调用，结合现场表现与成员共同完善转换条件和执行逻辑。\\
\textbf{林楗棋} & 动作组改进与转弯算法优化 & 调整与优化动作组及动作衔接，参与二维码转弯算法的实现与迭代；配合观察视觉判断与实际转向的对应关系，协同改进动作配置及算法细节，推进感知决策与机械执行的协调。\\
\textbf{曹淇涵} & 赛道搭建与动作辅助标定 & 参与围挡通道、转角及运行区域搭建和调整，协助开展动作标定与现场试验；配合核对机器人姿态、行走方向和转弯表现，将场地约束及运行现象反馈给标定与算法调试。\\
\textbf{刘小凡} & 算法协同构思、现场配合与报告撰写 & 参与二维码识别与转弯决策方案构思，协同开展动作标定、赛道搭建与现场调整；承担报告组织、技术机制梳理、图表制作和视频证据整理，将实现过程汇总为可复核的技术说明。\\\bottomrule
\end{longtable}
上述分工贯穿“方案讨论—动作标定—算法实现—场地验证—结果整理”全过程。动作和算法之间并非单向交付：转弯幅度影响视觉决策门限，围挡布局影响观测条件，现场反馈又推动动作组与程序调整。成员的主要方向彼此衔接，最终结果来自多轮共同讨论与交叉调试。

\subsection{其他小组成员与职责索引}
\begin{longtable}{p{25mm}p{34mm}p{81mm}}
\caption{第五大组其他小组成员与主要职责（按已提交资料）}\\\toprule
小组 & 成员/学号 & 主要职责\\\midrule\endfirsthead
\toprule 小组 & 成员/学号 & 主要职责\\\midrule\endhead
第二小组 & 陈凯华（学号待补） & AR尺寸、相机内参和TF标定与验证\\
第二小组 & \LuoMin{}（学号待补） & 跟随控制、速度门控、包迁移和参数调试\\
第二小组 & 刘彦廷（学号待补） & 命名空间、ID切换、队列场景与障碍测试\\
第二小组 & 许程奕（学号待补） & 实机试验、数据采集分析与报告撰写\\
第三小组 & 姓名/学号待补 & RGB-D感知、身份登记、跟随/导航及实验分别登记\\
方式四小组 & 徐杰鸿24363164 & 动态跟踪程序、避障优化、参数调试、试验与报告\\
方式四小组 & 王振宇24363096 & 标签配置安装、ID车辆对应与识别验证\\
方式四小组 & 林琢严24363054 & 辅助代码优化、场地布置\\
方式四小组 & 何爽24363028 & 辅助代码优化、场地布置\\\bottomrule
\end{longtable}
各方式章节保留对应组的详细分工；本章负责大组成员索引。最终人员、组号与学号应由各组统一核对，而非根据GitHub账户或画面车辆位置推测。

\subsection{跨组交付与共同责任}
第一小组向各跟随组提供领行方式、目标背贴位置及动作/任务状态变化；各跟随组反馈可见距离、最大可跟速度、转弯遮挡及停车条件。场地调整由全队同步，地图方式重新建图，视觉方式重新核对标记和身份模板。共同联调记录应包含每台车的控制器、直接目标、速度话题、停止方法和证据编号，由各模块负责人共同核对，报告整理由各组资料汇总完成。
